\unit{Exposición 236 -- Esquemas de Picard: propiedades generales}
\sid{236}
\src{221}
\seminares{14.º año, 1961/62, n.º 236\\Mayo de 1962}
\paperhead{TÉCNICA DE DESCENSO Y TEOREMAS DE EXISTENCIA EN GEOMETRÍA ALGEBRAICA\\
VI. ESQUEMAS DE PICARD: PROPIEDADES GENERALES}%
  {por Alexander GROTHENDIECK}

\fsection{0}{Complementos a la exposición anterior ([3], V)\textsuperscript{(1)}}

Se han producido algunos avances en las cuestiones de existencia de los
preesquemas de Picard planteadas en [3], V:

\medskip
\noindent a. (Mumford). No es cierto en general que, cuando
$f:X\longrightarrow S$ es un morfismo proyectivo separable (= plano con fibras
separables), exista el preesquema de Picard $\operatorname{Pic}_{X/S}$, ni
siquiera si las fibras de $f$ tienen dimensión 1 y $S$ es el espectro de un
anillo de valoración discreta completo. Se obtiene un ejemplo tomando
$S=\operatorname{Spec}\mathbf R[[t]]$ y tomando como $X$ el subesquema de
$\mathbf P^2_S$ definido por la ecuación (donde $x,y,z$ son variables
homogéneas)
\[
x^2+y^2=tz^2.
\]
Se trata de una cónica que degenera geométricamente en dos rectas que se
cortan, aunque la fibra especial sobre $\mathbf R$ es irreducible (está dada
por $x^2+y^2=0$ sobre $\mathbf R$). Se ve fácilmente que, tras la extensión
étale $S'\longrightarrow S$, con $S'=\operatorname{Spec}\mathbf C[[t]]$, existe
el preesquema de Picard de $X'/S'$ y admite una descripción explícita como
unión disjunta de copias de $\widetilde S$, donde $\widetilde S$ se obtiene de
$S$ duplicando el origen un número infinito de veces. Se comprueba fácilmente
que el dato de descenso sobre $\operatorname{Pic}_{X'/S'}$ para
$S'\longrightarrow S$ (dado aquí por la acción del grupo de Galois
$\mathbf Z/2\mathbf Z$ de $S'$ sobre $S$) no es efectivo: el grupo permuta
ciertos puntos duplicados, de modo que algunas órbitas no están contenidas en
ningún abierto afín. Mumford puede demostrar, sin embargo, que si
$f:X\longrightarrow S$ es un morfismo proyectivo separable tal que, para todo
$s\in S$, las componentes irreducibles de $X_s$ son geométricamente
irreducibles sobre $k(s)$, entonces existe $\operatorname{Pic}_{X/S}$. Su
demostración se apoya en un refinamiento de su teorema del cociente; cf. [7].
Obsérvese, por otra parte, que sin ninguna hipótesis sobre las componentes
irreducibles de las fibras $X_s$, sigue siendo posible que exista, no obstante,
el esquema $\operatorname{Pic}^{\tau}_{X/S}$ (que se introducirá más adelante).

\src{222}
\noindent b. (Murre). Sea $X$ un esquema propio sobre un cuerpo. Entonces
existe $\operatorname{Pic}_{X/k}$. La demostración sigue en parte la de
Chevalley [1] y utiliza de manera esencial la estructura de grupo del funtor de
Picard.

Para algunos complementos relativos a la teoría de los esquemas de Picard, en
particular en relación con los esquemas abelianos, será útil consultar [7]. Por
último, una laguna notable de la presente exposición es la ausencia de
«criterios de equivalencia» que permitan comparar el esquema de Picard de un
esquema proyectivo con los de sus secciones hiperplanas. Los teoremas clave
para desarrollar tales criterios se encuentran en [5], a los cuales hay que
añadir los teoremas de existencia para los esquemas de Picard.

\fsection{1}{Propiedades topológicas de los preesquemas conmutativos en grupos}

Sea primero $k$ un cuerpo y $G$ un preesquema en grupos sobre $k$. Como el
elemento neutro $e$, al ser racional sobre $k$, es necesariamente cerrado, se
deduce inmediatamente que la diagonal de $G\times_kG$ es cerrada y, por tanto,
que $G$ es separado: todo preesquema en grupos sobre un cuerpo es separado.
Denotemos por $G^\circ$ la componente conexa del elemento neutro $e$. Como $e$
es racional sobre $k$, $G^\circ$ es de hecho geométricamente conexa; es decir,
la formación de $G^\circ$ es compatible con la extensión del cuerpo de base.
Se deduce asimismo que $G^\circ$ es estable por multiplicación como conjunto y,
cuando $G$ es localmente noetheriano, de modo que $G^\circ$ es abierto, puede
considerarse un subgrupo abierto de $G$. Supongamos en lo sucesivo que $G$ es
localmente de tipo finito sobre $k$. Entonces $G^\circ$ es geométricamente
irreducible y de tipo finito sobre $k$. En efecto, se puede suponer que $k$ es
algebraicamente cerrado y después que $G$ es reducido (pues
$G_{\mathrm{red}}$ es entonces un subgrupo de $G$, ya que
$G_{\mathrm{red}}\times_kG_{\mathrm{red}}$ vuelve a ser reducido); por tanto,
es liso sobre $k$ en un abierto no vacío y, mediante traslación de ese
abierto, en todas partes. Así, $G$ es localmente irreducible, por lo que sus
componentes irreducibles coinciden con sus componentes conexas y $G^\circ$ es
irreducible. Sea $U$ un entorno afín de $e$ en $G^\circ$. Utilizando la
irreducibilidad de $G^\circ$, se ve inmediatamente que $U\cdot U=G^\circ$, lo
que demuestra que $G^\circ$ es cuasicompacto y, por tanto, de tipo finito sobre
$k$.

Para simplificar, supongamos que $G$ es conmutativo. Para todo entero $n>0$,
sea $G^{(n)}$ la imagen inversa de $G^\circ$ por el homomorfismo de potencia
$n$-ésima $\varphi_n$ de $G$; así, $G^{(n)}$ es un subgrupo abierto de $G$.
Pongamos
\[
G^\tau=\bigcup_nG^{(n)}.
\]

\src{223}
\[
G^\sigma=\bigcup_{(n,p)=1}G^{(n)},
\qquad
G^\rho=\bigcup_hG^{(p^h)},
\]
donde $p$ es el exponente característico del cuerpo $k$. Estos son subgrupos
abiertos de $G$ que satisfacen
\[
G^\sigma\cap G^\rho=G^\circ,
\qquad
G^\sigma G^\rho=G^\tau.
\]

\result{OBSERVACIÓN}{} Se puede construir el esquema en grupos cociente
$G/G^\circ=N$ (cf. [3], IV), y definir entonces $G^\tau$, $G^\sigma$ y
$G^\rho$ como las imágenes inversas en $G$, respectivamente, del subgrupo de
torsión de $N$, de la suma de sus componentes $q$-primarias para los primos
$q\neq p$ y de su componente $p$-primaria.\footnote{El texto impreso intercambia
las dos últimas correspondencias: asigna a $G^\sigma$ la componente
$p$-primaria y a $G^\rho$ su complemento primo con $p$. Las definiciones
precedentes de $G^\sigma$ y $G^\rho$ obligan al orden corregido que se da aquí.}
Obsérvese a este respecto que $N$ es un esquema en
grupos discreto separable sobre $k$. Así, tras elegir una clausura algebraica
$\overline k$ de $k$, con grupo de Galois $\pi$, se identifica con un grupo
discreto ordinario sobre el cual $\pi$ actúa por automorfismos. Esto da la
interpretación evidente de la construcción del subgrupo de torsión y de su
descomposición en componentes $q$-primarias. Cuando $G$ es el esquema de
Picard de un esquema propio $X$ sobre $k$, se puede llamar a $N$ el
\emph{esquema de Néron--Severi} (reducido) de $X$ sobre $k$. Cuando
$G^\circ_{\mathrm{red}}$ es un subesquema en grupos de $G^\circ$, lo que ocurre
en particular siempre que $k$ sea perfecto o $G^\circ$ sea propio sobre $k$
(por ejemplo, cuando $X$ sea geométricamente normal), conviene introducir
también el cociente $N'=G/G^\circ_{\mathrm{red}}$, que tiende a comportarse
mejor que $N$ por especialización, es decir, cuando $X$ varía en una familia
de esquemas algebraicos.

Sea ahora $S$ un preesquema localmente noetheriano y sea $G$ un preesquema en
grupos sobre $S$, localmente de tipo finito sobre $S$. No suponemos que $G$ sea
de tipo finito sobre $S$ ni separado sobre $S$. Pongamos
\[
G^\circ=\bigcup_{s\in S}(G_s)^\circ,
\]
y, cuando $G$ sea conmutativo,
\[
G^\tau=\bigcup_{s\in S}(G_s)^\tau,
\qquad
G^\sigma=\bigcup_{s\in S}(G_s)^\sigma,
\qquad
G^\rho=\bigcup_{s\in S}(G_s)^\rho.
\]

\src{224}
Estos son subconjuntos de $G$ estables por la multiplicación en $G$, lo que
evidentemente no implica que puedan ser definidos por subpreesquemas en grupos
de $G$. En particular, parece que en general no existe un subpreesquema en
grupos de $G$ cuyo conjunto subyacente sea $G^\circ$. Por supuesto, si uno de
estos conjuntos es abierto, entonces, con su estructura inducida, es un
subpreesquema abierto en grupos de $G$. Veremos que esto siempre ocurre para
$G^\tau$. Así, desde el punto de vista de los funtores representables, y en
particular desde el punto de vista de la especialización, la equivalencia
numérica se comporta de manera más satisfactoria que la equivalencia
algebraica. He aquí las principales propiedades generales de los conjuntos que
acabamos de definir:

\result{TEOREMA}{1.1} Los conjuntos $G^\circ$, $G^\tau$, $G^\sigma$ y $G^\rho$
son localmente constructibles. Además:

\noindent (i) $G^\circ$ es cuasicompacto sobre $S$. Cuando los $G_s^\circ$ son
propios y $G$ es separado sobre $S$, $G^\circ$ es propio sobre $S$ y, por
tanto, cerrado en $G$.

\noindent (ii) $G^\tau$ es abierto. Si $G^\circ$ es cerrado, también lo es
$G^\tau$.

\noindent (iii) Si $G^\circ$ es cerrado, también lo es $G^\sigma$, siempre que
se esté en característica igual, es decir, que todos los cuerpos residuales de
$S$ tengan la misma característica. Si $G^\circ$ es cerrado y
$G\longrightarrow S$ es universalmente abierto en los puntos de $G^\sigma$
(cf. Corolario 1.5 más adelante), entonces $G^\sigma\longrightarrow S$ es
universalmente abierto.

\noindent (iv) Si $G^\circ$ es cerrado, también lo es $G^\rho$. Supongamos que
se está en característica igual y que, para todo entero $n>0$ con $(n,p)=1$,
el homomorfismo de potencia $n$-ésima de $G$ es abierto. Entonces $G^\rho$ es
abierto.

Damos indicaciones para la demostración. La constructibilidad local de
$G^\circ$ y su cuasicompacidad sobre $S$ están contenidas en el lema siguiente:

\result{LEMA}{1.2} Supongamos $S\neq\varnothing$. Entonces existen un abierto
no vacío $U$ de $S$, un esquema en grupos $H$ de tipo finito sobre $U$ con
fibras conexas y un homomorfismo de esquemas en grupos
$H\longrightarrow G|U$ que es una inmersión abierta de imagen $G^\circ|U$.

Para demostrar el lema, se puede suponer que $S$ es irreducible y denotar por
$\eta$ su punto genérico. Tras el cambio de base
$S'=\operatorname{Spec}(\mathcal O_{S,\eta})\longrightarrow S$, se obtiene un
esquema en grupos $G'$ sobre el anillo local artiniano
$\mathcal O_{S,\eta}=A$, con un subesquema abierto en grupos $G'{}^\circ$ de
tipo finito sobre $A$, como se ha indicado antes. Por tanto, este procede de un
esquema en grupos $H$ de tipo finito sobre un entorno abierto $U$ de $\eta$, y
la inmersión canónica $G'{}^\circ\longrightarrow G'$ procede de una

\src{225}
inmersión abierta $H\longrightarrow G|U$, que es un homomorfismo de esquemas en
grupos después de hacer $U$ suficientemente pequeño. Si $U$ es suficientemente
pequeño, las fibras de $H$ son conexas y tienen todas la misma dimensión,
igual a la de las fibras de $G$. De ello se sigue que, para todo $s\in U$, la
imagen de $H_s$ en $G_s$ es exactamente $G_s^\circ$ (tras reducir de nuevo
$U$), lo que demuestra 1.2.

La segunda afirmación de (i) está contenida en el lema siguiente, aplicado a
un entorno abierto cuasicompacto de $G^\circ$ en $G$:

\result{LEMA}{1.3} Sea $X$ un preesquema separado de tipo finito sobre un
preesquema localmente noetheriano $S$, sea $g$ una sección de $X$ sobre $S$ y
sea $X^\circ$ la unión, en los $X_s$, de las componentes conexas que contienen
$g(s)$. Si $s\in S$ y $X_s^\circ$ es propio sobre $k(s)$, existe un entorno
abierto $U$ de $s$ tal que $X^\circ|U$ es propio sobre $U$ y, por tanto,
cerrado en $X|U$.

Por descenso fielmente plano sobre la base, se reduce al caso en que $S$ es el
espectro de un anillo local completo y $s$ es su punto cerrado. Aplicando [2],
III, 5.5.1, se ve que $X$ es la unión disjunta de dos abiertos $X'$ y $X''$, el
primero propio sobre $S$ y tal que $X'_s=X_s^\circ$. Esto reduce al caso
$X=X'$, es decir, a $X$ propio sobre $S$. Se utiliza entonces un argumento
estándar basado en el criterio valuativo de ser propio para un subconjunto
(omitido en [2], capítulo II).

Demostremos ahora que $G^\tau$ es abierto. Equivalentemente, puesto que la
formación de $G^\tau$ (así como la de $G^\circ$, $G^\sigma$ y $G^\rho$) conmuta
con la extensión de la base, para toda sección $g$ de $G$ sobre $S$, el
conjunto $g^{-1}(G^\tau)$ es abierto. Esto significa dos cosas:

\medskip
\noindent a. Si $y\in S$ y $g(y)\in G_y^\tau$, entonces
$g(y')\in G_{y'}^\tau$ para todo $y'\in\overline{\{y\}}$ suficientemente
próximo a $y$.

\medskip
\noindent b. Si $y'\in S$ y $g(y')\in G_{y'}^\tau$, entonces
$g(y)\in G_y^\tau$ para toda generalización $y$ de $y'$.

Para (a), existe un entero $n>0$ tal que $g^n(y)\in G_y^\circ$. Como
$G^\circ$ es constructible, $(g^n)^{-1}(G^\circ)$ es constructible y, por
tanto, $g^n(y')\in G_{y'}^\circ$ para
$y'\in\overline{\{y\}}$ suficientemente próximo a $y$. Para (b), las
potencias $g(y')^n=g^n(y')$ permanecen en un abierto cuasicompacto de
$G_{y'}$, pues están contenidas en un número finito de clases módulo
$G_{y'}^\circ$. Por consiguiente, existe un abierto cuasicompacto $U$ de $G$
que contiene los $g^n(y')$ y, por tanto, también sus generalizaciones
$g^n(y)$. Así, las potencias de $g(y)$ permanecen en un abierto cuasicompacto
de $G_y$, lo que implica fácilmente que $g(y)\in G_y^\tau$.

\src{226}
Supongamos que $G^\circ$ es cerrado; demostremos lo mismo para $G^\tau$. Como
ya se sabe que $G^\tau$ es abierto y, por tanto, localmente constructible,
basta demostrar que es estable por especialización. Esto se sigue de que es
una unión de conjuntos cerrados, a saber, las imágenes inversas del conjunto
cerrado $G^\circ$ por los homomorfismos de potencia $n$-ésima $\varphi_n$.

El mismo argumento demuestra que $G^\sigma$ y $G^\rho$ son cerrados si
$G^\circ$ es cerrado (en el primer caso, bajo la condición de característica
igual), una vez demostrado que $G^\sigma$ y $G^\rho$ son localmente
constructibles. Para $x\in G^\tau$, sea $\nu(x)$ el menor entero $n>0$ tal que
el homomorfismo de potencia $n$-ésima $\varphi_n$ envía $x$ a $G^\circ$.
Entonces $G^\sigma$ (respectivamente, $G^\rho$) está formado por los
$x\in G$\footnote{El texto impreso dice «$X$», que no está definido en este
pasaje; los cuatro loci considerados son subconjuntos de $G$.} para los cuales
$\nu(x)$ es primo con $p$ (respectivamente, una potencia de $p$). La afirmación
de constructibilidad se deduce ahora del enunciado más preciso:

\result{LEMA}{1.4} La función $\nu$ sobre $G^\tau$ es localmente
constructible.

En efecto, para todo entero $n>0$, el conjunto de los $x\in G^\tau$ tales que
$\nu(x)=n$ es localmente constructible: es la diferencia entre
$\varphi_n^{-1}(G^\circ)$ y la unión de los
$\varphi_d^{-1}(G^\circ)$, donde $d$ recorre los divisores propios de $n$.
Como $G^\circ$ es localmente constructible, también lo son los
$\varphi_d^{-1}(G^\circ)$ y, por consiguiente, también lo es la diferencia
anterior.

Supongamos ahora que $G^\circ$ es cerrado y que $G\longrightarrow S$ es
universalmente abierto en los puntos de $G^\sigma$. Demostremos que
$G^\sigma\longrightarrow S$ es abierto, es decir, que envía un entorno en
$G^\sigma$ de $x\in G^\sigma$ sobre un entorno de $y=f(x)$. Como $G^\sigma$ es
localmente constructible, basta demostrar que para toda generalización $y'$ de
$y$ existe una generalización $x'$ de $x$ en $G^\sigma$ situada sobre $y'$.
Por cambio de base, se reduce al caso en que $S$ es el espectro de un anillo de
valoración discreta y $y,y'$ son, respectivamente, sus puntos cerrado y
genérico. Como $G\longrightarrow S$ es abierto en $x$, existe una
generalización $x_1$ de $x$ en $G$ sobre $y'$. Tras otro cambio de base, se
puede suponer además que existe una sección $g$ de $G$ sobre $S$ con
$x=g(y)$. Si $k(y')$ tiene característica cero, tómese cualquier
generalización $x'$ de $x$ en $G$ sobre $y'$. Pertenece a $G^\tau$ porque
$G^\tau$ es abierto y, por tanto, a $G^\sigma$ porque
$G_{y'}^\sigma=G_{y'}^\tau$. Si $k(y')$ tiene característica $p>0$, escribamos
\[
\nu(g(y'))=p^hm\qquad\text{con }(m,p)=1;
\]

\src{227}
elijamos enteros $a$ y $b$ tales que $ap^h+bm=1$, y pongamos
\[
g_1=g^{ap^h},\qquad g_2=g^{bm},\qquad\text{de modo que}\qquad g=g_1g_2.
\]
Por construcción, $g_1(y')\in G^\sigma$ y $g_2(y')\in G^\rho$. Como
$G^\circ$ es cerrado, se sigue que $g_1(y)\in G^\sigma$ y
$g_2(y)\in G^\rho$. Pero
\[
g(y)=g_1(y)g_2(y)\in G^\sigma,
\]
por lo que $g_2(y)\in G^\sigma$ y, por consiguiente,
\[
g_2(y)\in G_y^\sigma\cap G_y^\rho=G_y^\circ.
\]
La hipótesis, junto con el hecho de que $S$ es el espectro de un anillo de
valoración discreta, implica que $G-(G_y-G_y^\circ)$ es un abierto de $G$ en
el que $G\longrightarrow S$ induce un morfismo abierto. Así, por todo punto de
$G_y^\circ$ pasa una «cuasisección». Después de efectuar una extensión
adicional de la base $S$, se puede suponer, por consiguiente, que existe una
sección $g'_2$ de $G^\circ$ sobre $S$ tal que $g'_2(y)=g_2(y)$. Pongamos
$g'=g_1g'_2$. Entonces, por construcción, $g'(y)=g(y)=x$, y
\[
g'(y')=g_1(y')g'_2(y')\in G^\sigma.
\]
Por tanto, $g'(y')=x'$\footnote{El texto impreso dice «$g(y')=x'$». La sección
construida en la frase anterior, cuya fibra genérica se calcula, es $g'$.} es
una generalización de $x$ en $G^\sigma$ situada sobre $y'$. Esto demuestra
(iii).

Queda por demostrar la última afirmación, (iv). Debemos probar que, si
$x\in G^\rho$, toda generalización $x'$ de $x$ pertenece a $G^\rho$. Tomando
imágenes por $\varphi_{p^h}$ para un $h$ adecuado\footnote{El texto impreso
dice $\varphi_{n^h}$. Para $x\in G^\rho$, la definición de la p.~223 da una
potencia del exponente característico $p$ que envía $x$ a $G^\circ$.}, se
puede suponer que, de hecho, $x\in G^\circ$. Para todo entero $n$ primo con la
característica, $x$ pertenece entonces a la imagen de $\varphi_n$ (pues la
aplicación de potencia $n$-ésima en un grupo conexo de tipo finito sobre un
cuerpo de característica prima con $n$ es sobreyectiva). Como $\varphi_n$ es
abierto, $x'$ pertenece también a la imagen de $\varphi_n$. Más precisamente,
sea $U$ un abierto cuasicompacto de $G^\tau$ que contenga $G_y^\circ$.
Entonces $x'\in\varphi_n(U)$ para todo $n$ primo con $p$. Tomando por $n$ un
múltiplo común de las partes primas con $p$ de los $\nu(z)$, para $z\in U$, se
obtiene $x'\in G^\rho$.

La afirmación (1.1), (iii), se completa mediante el corolario siguiente:

\result{COROLARIO}{1.5} Sea $n>1$ un entero tal que el homomorfismo de potencia
$n$-ésima $\varphi_n:G\longrightarrow G$ sea universalmente abierto (por
ejemplo, étale), pongamos
\[
G^{(n)}=\varphi_n^{-1}(G^\circ),
\]
y supongamos que, en cada fibra geométrica $G_{\bar s}$, el subconjunto
$\bigcup_{h\geq0}|{}_{n^h}G_{\bar s}|$ es denso en
$\bigcup_{h\geq0}G_{\bar s}^{(n^h)}$. Entonces

\src{228}
$G\longrightarrow S$ es universalmente abierto en los puntos de los
$G^{(n^h)}$. En particular, si los $G_s^\circ$ no contienen ninguna componente
aditiva y, para todo entero $n>1$, el homomorfismo $\varphi_n$ es
universalmente abierto en los puntos de característica residual prima con
$n$\footnote{El texto impreso dice «prima con $p$». En este enunciado
cuantificado por $n$, el criterio para que un morfismo sea étale y el argumento que utiliza los
núcleos $ {}_nG$ requieren característica residual prima con $n$.}, entonces
$G\longrightarrow S$ es universalmente abierto en los puntos de $G^\sigma$.
Por tanto, según (1.1), (iii), si $G^\circ$ es cerrado,
$G^\sigma\longrightarrow S$ es universalmente abierto. En estos casos,
$s\longmapsto\dim G_s$ es localmente constante sobre $S$.

En efecto, la hipótesis implica que el núcleo $ {}_nG$ de $\varphi_n$ es
universalmente abierto sobre $S$. Por tanto, $G\longrightarrow S$ es
universalmente abierto en los puntos de $ {}_nG$ y, reemplazando $n$ por
$n^h$, en los puntos de $ {}_{n^h}G$. La hipótesis de densidad dice exactamente que, fibra a fibra, los puntos cuyo
orden es una potencia de $n$ son densos en la unión de los $G^{(n^h)}$. Para
la afirmación particular, la ausencia de componentes aditivas implica esta
densidad para los enteros $n$ invertibles en el campo residual. Se deduce fácilmente que
$G\longrightarrow S$ es universalmente abierto en cada punto de esta unión,
en particular a lo largo de la sección neutra, y, por tanto, que
$s\longmapsto\dim G_s$ es localmente constante.

\result{OBSERVACIÓN}{1.6} Recuérdese que un morfismo $X\longrightarrow S$ se
llama universalmente abierto si envía todo abierto a un abierto y conserva
esta admirable propiedad después de cualquier cambio de base
$S'\longrightarrow S$. Esto también significa (cuando $S$ es localmente
noetheriano y $X\longrightarrow S$ es localmente de tipo finito) que toda
componente irreducible de $X$ domina $S$ y que esta propiedad se conserva
después de cualquier cambio de base $S'\longrightarrow S$. En ambas
afirmaciones basta (bajo las hipótesis de finitud anteriores) comprobar los
cambios de base $S'\longrightarrow S$ en los que $S'$ es el espectro de un
anillo de valoración discreta (completo y con cuerpo residual algebraicamente
cerrado, si se desea~...). La definición se extiende de manera evidente a un
subconjunto $Z$ de $X$ (como el subconjunto $G^\sigma$ de $G$). Es un fenómeno
perfectamente normal, incluso para un morfismo proyectivo liso
$X\longrightarrow S$ con fibras geométricamente conexas (por ejemplo, el
producto fibrado de la familia modular de curvas elípticas consigo misma sobre
$S=\operatorname{Spec}\mathbf C[j]$), que $\operatorname{Pic}_{X/S}$ no sea
universalmente abierto sobre $S$: una componente irreducible de
$\operatorname{Pic}_{X/S}$ puede estar situada enteramente sobre un único
punto de $S$. Esto está relacionado con el hecho de que el rango del grupo de
Néron--Severi de las fibras de $X/S$ puede aumentar bruscamente (fenómenos de
«multiplicación compleja»). Por otra parte, (1.5) nos asegura que, en casos
favorables, $\operatorname{Pic}_{X/S}^\sigma$ (y parece que muy a menudo incluso
$\operatorname{Pic}_{X/S}^\tau$) es universalmente abierto sobre $S$.

\src{229}
Por último, he aquí un caso útil en el que, excepcionalmente, $G^\circ$
consiente en ser abierto:

\result{COROLARIO}{1.7} Supongamos que $G\longrightarrow S$ es universalmente
abierto en los puntos de $G^\circ$ (cf. (1.5)) y que las fibras $G_s$ son
separables y, por tanto, lisas sobre $k(s)$ (esta última condición se satisface
automáticamente en característica residual cero, por un teorema de Cartier).
Entonces $G^\circ$ es abierto en $G$. Si, además, $S$ es reducido, entonces
$G^\circ$ es liso (en particular, plano) sobre $S$.

La primera afirmación puede precisarse del modo siguiente. Si, para un
$s\in S$ dado, $G\longrightarrow S$ es universalmente abierto en los puntos de
$G_s^\circ$ y $G_s^\circ$ es separable sobre $k(s)$, entonces $G^\circ$ es un
entorno de $G_s^\circ$ (y la hipótesis hecha en $s$ seguirá siendo válida en
los puntos próximos). Este enunciado, independiente de toda estructura de
grupo, se demostrará en [2], IV, §7. La última afirmación de (1.7), también
independiente de toda estructura de grupo y de toda cuestión relativa a una
componente conexa, es un caso particular de un criterio de planitud dado en
[2], IV, §5, que implica de manera más general:

\result{COROLARIO}{1.8} Sea $U$ un abierto de una fibra $G_s$ tal que
$G\longrightarrow S$ sea universalmente abierto en los puntos de $U$ (cf.
(1.5)). Si $G_s$ es separable sobre $k(s)$ y $S$ es reducido en $s$, entonces
$G$ es plano sobre $S$ en los puntos de $U$ (y, por tanto, en este caso, liso
sobre $S$ en los puntos de $U$).

\result{OBSERVACIONES}{1.9} No sé si, cuando $G$ es separado sobre $S$,
$G^\circ$ o $G^\tau$ es siempre cerrado en $G$; parece poco probable. En todo
caso, hay contraejemplos evidentes si se suprime la hipótesis de separación.
Por ejemplo, tómese la recta afín con el origen duplicado una infinidad
numerable de veces. Esto da un esquema en grupos $G$ sobre la recta afín cuyas
fibras se reducen todas al grupo neutro salvo una, cuya fibra es $\mathbf Z$.
Este esquema en grupos es un subesquema abierto en grupos, a saber, la clausura
de la sección neutra, del preesquema de Picard del $S$-esquema $X$
correspondiente a una familia de cónicas que degenera en dos rectas que se
cortan.

Además, incluso partiendo de un esquema en grupos finito y plano $G$ sobre el
espectro $S$ de un anillo de valoración discreta $V$, de modo que $G^\circ$ se
reduce a la sección neutra y es, por tanto, cerrado, varias conclusiones pasan
a ser falsas si se suprimen ciertas hipótesis. Supongamos que $V$ tiene
característica igual $p>0$ y sea $G$ el núcleo del homomorfismo
$\mathbf G_a\longrightarrow\mathbf G_a$ definido por el homomorfismo de
funtores
\[
f\longmapsto f^p-tf,
\]
donde $t$ es un uniformizante de $V$.

\src{230}
(Recuérdese que, por definición, el «grupo aditivo» $\mathbf G_a$ sobre $S$
representa el funtor $S'\longmapsto\Gamma(S',\mathcal O_{S'})$.) Así, como
$S$-esquema,
\[
G=\operatorname{Spec}V[x]/(x^p-tx)
\]
es la unión de $p$ «rectas» que se cortan, y representa un grupo
$\mathbf Z/p\mathbf Z$ que degenera en un grupo infinitesimal de tipo aditivo.
En este ejemplo, $G^\circ=G^\sigma$ (y es una de las $p$ rectas), por lo que no
es abierto en $G$. En característica mixta de característica residual $p>0$,
partamos del esquema en grupos $H=\boldsymbol\mu_p$, núcleo de la aplicación de
potencia $p$-ésima en $\mathbf G_m$. Así,
\[
H=\operatorname{Spec}V[x]/(x^p-1);
\]
Este es de nuevo la unión de $p$ rectas que se cortan, y representa un grupo
$\mathbf Z/p\mathbf Z$ (en característica cero) que degenera en un grupo
infinitesimal de tipo multiplicativo en característica $p$. Sea $H'$ el
«esquema en grupos constante» definido por el grupo finito ordinario
$\mathbf Z/p\mathbf Z$, es decir, la unión disjunta de $\mathbf Z/p\mathbf Z$
copias de $S$, o equivalentemente
$\operatorname{Spec}V^{\mathbf Z/p\mathbf Z}$. Entonces $G=H\times_S H'$
describe un grupo $(\mathbf Z/p\mathbf Z)^2$ en característica cero, que
degenera en característica $p$ en el producto de un grupo infinitesimal por
$\mathbf Z/p\mathbf Z$. Aquí $G^\rho$ es la unión de la sección neutra y la
fibra especial, de modo que $G^\rho$ no es abierto, mientras que $G^\sigma$ es
la unión de la sección neutra y la fibra genérica, de modo que no es cerrado,
contrariamente a lo afirmado en los casos de característica igual de (1.1),
(iii) y (iv). Se trata, naturalmente, de fenómenos relacionados con la
característica $p>0$. Los resultados anteriores dan:

\result{COROLARIO}{1.10} Supongamos que todas las características residuales de
$S$ son cero, de modo que $G^\tau=G^\sigma$ y $G^\circ=G^\rho$. Entonces
$G^\tau=G^\sigma$ es abierto; es incluso abierto y cerrado si $G$ es separado
sobre $S$ y los $G_s^\circ$ son propios. Bajo la misma hipótesis,
$G^\circ=G^\rho$ es propio sobre $S$ y, por tanto, cerrado en $G$. Supongamos,
por último, que para todo entero $n>1$ el homomorfismo de potencia $n$-ésima de
$G$ es universalmente abierto. Entonces $G^\circ$ es abierto; y si, además,
los $G_s^\circ$ no tienen componente aditiva (por ejemplo, si los
$G_s^\circ$ son propios, como antes), entonces $G^\tau\longrightarrow S$ es
universalmente abierto, e incluso liso si $S$ es reducido.

Consignemos finalmente el siguiente resultado fácil:

\result{PROPOSICIÓN}{1.11} Existe un abierto $U$ de $S$ tal que el conjunto
$G'$ de los puntos de $G$ en los que $G$ es liso (respectivamente, plano) sobre
$S$ es el conjunto subyacente de un subesquema abierto en grupos inducido por
$G|U$. Además, toda sección de $G$ sobre $U$ es una sección de $G'$ sobre $U$.

\src{231}
\result{COROLARIO}{1.12} Si $G$ es liso (respectivamente, plano) sobre $S$ en
los puntos de la sección neutra, también lo es en los puntos de toda sección de
$G$ sobre $S$ y en todo punto de $G^\circ$. Si, además, para todo entero $n>0$
el homomorfismo de potencia $n$-ésima
$\varphi_n:G\longrightarrow G$ es étale en los puntos de característica prima
con $n$, entonces $G$ es liso (respectivamente, plano) sobre $S$ en todo punto
de $G^\sigma$.

\fsection{2}{Aplicación a las propiedades locales de los esquemas de Picard}

\result{TEOREMA}{2.1}

\noindent(i) Sea $f:X\longrightarrow S$ un morfismo propio y liso, y
supongamos que existe $\operatorname{Pic}_{X/S}$ (por ejemplo, supongamos que
$f$ es proyectivo). Entonces $\operatorname{Pic}_{X/S}$ es separado sobre $S$
y todo subconjunto cerrado $Z$ de $\operatorname{Pic}_{X/S}$ que sea de tipo
finito sobre $S$ es propio sobre $S$.

\noindent(ii) Sea $X$ un preesquema propio geométricamente normal sobre un
cuerpo $k$. Entonces $\operatorname{Pic}_{X/k}^\circ$\footnote{El texto
impreso dice $\operatorname{Pic}_{X/S}^\circ$, aunque $S$ no aparece en este
enunciado sobre el cuerpo $k$.} es propio sobre $k$.

\begin{proof}
(i) Los criterios valuativos ([2], II, §7) nos reducen a demostrar lo
siguiente. Sea $S$ el espectro de un anillo de valoración discreta completo y
sea $U$ el abierto formado por el punto genérico de $S$. Toda sección racional
de $\operatorname{Pic}_{X/S}$ sobre $S$, es decir, toda sección sobre $U$, se
prolonga de manera única a una sección sobre $S$. Por la definición de
$\operatorname{Pic}_{X/S}$, esto equivale al enunciado siguiente: para todo
Módulo invertible $\mathcal L$ sobre $V=f^{-1}(U)$ existe un Módulo invertible sobre $X$
que prolonga $\mathcal L$, único salvo isomorfismo. Esto se deduce fácilmente
de la descripción de los Módulos invertibles sobre $V$ y, respectivamente, sobre
$X$, en términos de clases de divisores de Cartier. En efecto, los anillos
locales de $X$ son regulares ($X$ es liso sobre el esquema regular $S$) y, por
tanto, factoriales por el teorema de Auslander; así, todo divisor sobre
$X$\footnote{El texto impreso dice «sobre $S$». La factorialidad invocada es
la de los anillos locales de $X$.} es un divisor de Cartier. Todo divisor sobre
$V$ se prolonga a un divisor sobre $X$ tomando su clausura.
\end{proof}

\result{OBSERVACIÓN}{2.2} La demostración sigue siendo válida si se supone
solamente que $f$ es plano y que sus fibras $X_s$ son intersecciones completas
locales y lisas sobre $k(s)$ en codimensión $\leq2$. Esto se deduce del teorema
demostrado en [5]: un anillo local noetheriano de intersección completa que sea
regular en codimensión $\leq3$ es factorial («conjetura de Samuel»). Obsérvese
que el

\src{232}
resultado pasa a ser falso si se reemplaza «codimensión $\leq2$» por
«codimensión $\leq1$», es decir, por la hipótesis «normal», como muestra una
familia de cuádricas no singulares que degenera en un cono cuadrático.

(ii) El lema de Chow nos reduce al caso en que $X$ es proyectivo y, por tanto,
está inmerso en $\mathbf P_k^n$; podemos suponer que $X$ es conexo. Si
$\dim X=1$, entonces $X$ es liso sobre $k$ y se aplica (i). Si $\dim X\geq2$,
una variante de los «criterios de equivalencia» conocidos muestra que existe
un número finito de curvas lisas $Y_i$ contenidas en $X$ (obtenidas como
intersecciones de $X$ con subespacios lineales adecuados de $\mathbf P_k^n$),
tales que
\[
\operatorname{Pic}_{X/k}^\tau\longrightarrow
\prod_i\operatorname{Pic}_{Y_i/k}^\tau
\]
es un monomorfismo y, por tanto, induce un monomorfismo sobre las componentes
conexas. La componente conexa del segundo miembro es propia sobre $k$ por el
resultado anterior, y el homomorfismo de esquemas en grupos considerado es
necesariamente una inmersión cerrada. Se sigue que
$\operatorname{Pic}_{X/k}^\circ$ es también propio sobre $k$. Los delicados
criterios de equivalencia pueden evitarse utilizando la estructura de los
grupos algebraicos conmutativos sobre un cuerpo algebraicamente cerrado
(debida a Chevalley--Borel). Se reduce a demostrar que todo morfismo de la
recta afín privada del origen\footnote{El texto impreso dice «de la recta
afín». La fe de erratas vinculada, p.~303, prescribe «recta afín privada del
origen».} a $\operatorname{Pic}_{X/k}^\tau$ es constante. Esto equivale a decir
que todo Módulo invertible sobre $X[t,t^{-1}]$\footnote{El texto impreso dice
$X[t]$; la fe de erratas vinculada, p.~303, prescribe $X[t,t^{-1}]$.} procede
de un Módulo invertible sobre $X$. Este conocido resultado elemental ni siquiera
utiliza que $X$ sea propio sobre $k$ (la normalidad de $X$ reduce
inmediatamente al caso en que $X$ es regular).

\result{COROLARIO}{2.3} Sea $f:X\longrightarrow S$ un morfismo propio normal
(es decir, plano con fibras geométricamente normales), y supongamos que existe
$\operatorname{Pic}_{X/S}$. Entonces $\operatorname{Pic}_{X/S}^\circ$ es propio
sobre $S$ y, por tanto, cerrado en $\operatorname{Pic}_{X/S}$. Además,
$\operatorname{Pic}_{X/S}^\tau$ y $\operatorname{Pic}_{X/S}^\rho$ son también
cerrados, al igual que $\operatorname{Pic}_{X/S}^\sigma$ en «característica
igual».

Aplíquense (1.1) y (2.1), (ii).

\result{COROLARIO}{2.4} Sea $f:X\longrightarrow S$ un morfismo propio y liso
tal que exista $\operatorname{Pic}_{X/S}$ y que este sea una unión disjunta de
esquemas $P^{(i)}$ de tipo finito sobre $S$ (cf. [3], V, 4.1). Entonces cada
$P^{(i)}$ es propio sobre $S$.

Esto se deduce de (2.1), (i). Como se señaló en (2.2), el resultado puede
generalizarse imponiendo hipótesis menos restrictivas sobre las fibras de $f$,
pero pasa a ser falso si se supone meramente que $f$ es normal. En ese caso,
no sé si $\operatorname{Pic}_{X/S}^\tau$ es, no obstante, propio sobre $S$,
incluso suponiendo que sea de tipo finito sobre $S$.

\src{233}
\result{TEOREMA}{2.5} Sea $f:X\longrightarrow S$ un morfismo propio y plano tal
que exista $\operatorname{Pic}_{X/S}$, y para cada entero $n$ sea $\varphi_n$
el homomorfismo de potencia $n$-ésima de este preesquema en grupos. Entonces
$\varphi_n$ es étale en todo punto
$x\in\operatorname{Pic}_{X/S}$\footnote{El texto impreso dice «$x\in X$». La
aplicación $\varphi_n$ es el endomorfismo de potencia del preesquema de
Picard.} cuya característica residual sea prima con $n$.

Por la caracterización infinitesimal de los morfismos étales, esta afirmación
equivale a la siguiente:

\result{LEMA}{2.6} Supongamos que $S$ es el espectro de un anillo artiniano
local $A$ cuyo ideal maximal $\mathfrak m$ tiene potencia $(\nu+1)$-ésima nula,
y pongamos
\[
A_{\nu-1}=A/\mathfrak m^\nu,
\qquad X_{\nu-1}=X\otimes_A A_{\nu-1}.
\]
Sea $\mathcal L$ un Módulo invertible sobre $X$ y sea
$\mathcal L'_{\nu-1}$ un Módulo invertible sobre $X_{\nu-1}$ cuya potencia
tensorial $n$-ésima sea isomorfa a
\[
\mathcal L_{\nu-1}=\mathcal L\otimes_A A_{\nu-1}.
\]
Entonces existe un Módulo invertible $\mathcal L'$ sobre $X$ cuya potencia
tensorial $n$-ésima es isomorfa a $\mathcal L$, siempre que $n$ sea primo con
la característica residual de $k=k(A)$.

Demostremos el lema. Pongamos
\[
V=\mathfrak m^\nu=\mathfrak m^\nu/\mathfrak m^{\nu+1};
\]
este es un espacio vectorial sobre $k=k(A)$. Intentemos primero prolongar
$\mathcal L'_{\nu-1}$ a un Módulo invertible arbitrario $\mathcal L'$ sobre $X$.
La obstrucción pertenece a
\[
H^2(X_0,\mathcal O_{X_0})\otimes_k V.
\]
La hipótesis sobre $\mathcal L'_{\nu-1}$ y el hecho de que
$\mathcal L_{\nu-1}$ pueda prolongarse muestran que $n$ veces esta obstrucción
es cero; la propia obstrucción es, por tanto, nula porque $n$ es primo con la
característica. La ambigüedad en la prolongación elegida pertenece a
\[
H^1(X_0,\mathcal O_{X_0})\otimes_k V,
\]
y la diferencia $\xi$ entre $\mathcal L'^{\otimes n}$ y $\mathcal L$ pertenece
al mismo Módulo. Corregir $\mathcal L'$ para anular esta diferencia equivale a
encontrar un $\eta$ en ese Módulo tal que $n\eta=\xi$. Esto vuelve a ser
posible porque $n$ es primo con la característica.

\result{COROLARIO}{} Bajo las hipótesis de (2.5), supongamos además que los
esquemas de Picard de las fibras $X_s$ no contienen ninguna componente aditiva
(por ejemplo, que los $X_s$ son geométricamente normales; cf. (2.1), (ii)).
Entonces $\operatorname{Pic}_{X/S}\longrightarrow S$ es universalmente abierto
en los puntos de $\operatorname{Pic}_{X/S}^\circ$. Si
$\operatorname{Pic}_{X/S}^\circ$ es cerrado (por ejemplo, si los $X_s$ son
geométricamente normales; cf. (2.3)), entonces
$\operatorname{Pic}_{X/S}^\sigma$ es universalmente abierto sobre $S$. Por
último, en característica igual,
$\operatorname{Pic}_{X/S}^\rho\longrightarrow S$ es universalmente abierto.

Aplíquense (1.5) y (1.1).

\src{234}
\result{COROLARIO}{2.7} Sea $f:X\longrightarrow S$\footnote{El texto impreso
dice «$f:X\longrightarrow Y$», aunque tanto el preesquema de Picard como la
función siguiente tienen base $S$.} un morfismo propio y plano tal que exista
$\operatorname{Pic}_{X/S}$. Entonces la función
\[
s\longmapsto\dim\operatorname{Pic}_{X_s/k(s)}
\]
sobre $S$ es semicontinua superiormente (es decir, puede saltar hacia arriba,
pero no hacia abajo), e incluso continua (es decir, localmente constante) si
los $\operatorname{Pic}_{X_s/k(s)}$ no contienen ninguna componente aditiva.

La primera afirmación es trivial, o casi, para todo preesquema en grupos
localmente de tipo finito sobre una base localmente noetheriana, pues basta
mirar a lo largo de la sección neutra. La segunda afirmación se deduce de
(2.5).

\result{OBSERVACIÓN}{2.8} Sean $s,s'\in S$, con $s$ una especialización de
$s'$. Entonces (2.7) equivale a una desigualdad (respectivamente, una igualdad)
entre las dimensiones de los esquemas de Picard de $X_{s'}$ y de su
«especialización» $X_s$. Incluso antes de que se construyeran los esquemas de
Picard, Serre había observado que, para un morfismo liso
$f:X\longrightarrow S$, la invariancia de la dimensión de los Picard de los
$X_s$ es una consecuencia formal de la teoría de la especialización del grupo
fundamental ([4], X) y de las relaciones clásicas de Kummer entre los puntos de
orden finito de la variedad de Picard y el grupo fundamental abelianizado
([4], XI). En general, sean $\alpha,\mu,\lambda$ las dimensiones de las partes
abeliana, multiplicativa y aditiva de
$\operatorname{Pic}_{X_s/k(s)}^\circ$, y definamos de modo análogo
$\alpha',\mu',\lambda'$. Las relaciones conocidas se expresan mediante las
desigualdades siguientes:
\[
\tag{*}\alpha+\mu+\lambda\geq\alpha'+\mu'+\lambda'
\]
(válida siempre que exista $\operatorname{Pic}_{X/S}$ y, por tanto,
probablemente en todos los casos). Cuando $\lambda=\lambda'=0$, esta
desigualdad se convierte en una igualdad, válida bajo las mismas hipótesis de
existencia:
\[
\alpha+\mu=\alpha'+\mu';
\]
por otra parte,
\[
\tag{**}2\alpha+\mu\leq2\alpha'+\mu'
\]

\src{235}
por el argumento de Serre si los $X_s$ son separables (sin suponer siquiera la
existencia de $\operatorname{Pic}_{X/S}^\tau$), o si los
$ {}_n\operatorname{Pic}_{X/S}$ (los núcleos de $\varphi_n$ en el funtor de
Picard) son separados sobre el abierto de $S$ donde $n$ es invertible (pues son étales sobre ese abierto por (2.5)). Se
inclina uno a conjeturar que (*) es siempre una igualdad, o al menos cuando los
$X_s$ son separables, y que
\[
\tag{***}\alpha\leq\alpha',\qquad\lambda\geq\lambda'
\]
se verifica siempre que exista un preesquema en grupos, localmente de tipo
finito sobre un $S$ localmente noetheriano, cuyas fibras tengan dimensión
constante (véase (1.3) para un resultado positivo en esta dirección).

\result{OBSERVACIÓN}{2.9} En todos los casos conocidos,
$\operatorname{Pic}_{X/S}^\tau$ es universalmente abierto sobre $S$, pero
probablemente no conviene dejarse llevar por ilusiones excesivas, ni siquiera
cuando $f:X\longrightarrow S$ es liso. En todo caso, Mumford ha construido un
ejemplo (ciertamente con $S$ no reducido, de hecho el espectro de un anillo
artiniano) en el que $\operatorname{Pic}_{X/S}^\tau$ no es plano sobre $S$,
haciendo variar infinitesimalmente la superficie de Igusa. El punto considerado
por Mumford pertenece a $\operatorname{Pic}_{X/S}^\rho$, y sigue siendo posible
que, para un $f:X\longrightarrow S$ liso, $\operatorname{Pic}_{X/S}^\tau$ sea
plano sobre $S$ en los puntos de $\operatorname{Pic}_{X/S}^\sigma$. El ponente
duda, sin embargo, que esto sea siempre así, incluso si se restringe a los
puntos de $\operatorname{Pic}_{X/S}^\circ$ y se supone que $S$ es el espectro
de un anillo de valoración discreta. La cuestión está también relacionada con
el estudio de los puntos fijos de un esquema abeliano bajo un grupo finito de
automorfismos, situación para la que parecen faltar ejemplos.

Parecería que, incluso restringiéndose a un morfismo proyectivo liso
$f:X\longrightarrow S$, los resultados de regularidad local para
$\operatorname{Pic}_{X/S}$ enunciados en este parágrafo, junto con las
conjeturas de (2.8), agotan casi todo lo que puede decirse sin hipótesis más
específicas sobre la naturaleza de las fibras de $f$. Recuérdese, sin embargo,
que si las fibras geométricas de $\operatorname{Pic}_{X/S}$ son reducidas y no
tienen componente aditiva, entonces (1.8) y (2.5) implican que
$\operatorname{Pic}_{X/S}$ es liso sobre $S$ en los puntos de
$\operatorname{Pic}_{X/S}^\sigma$ cuando $S$ es reducido. El mismo resultado se
verifica, sin ninguna hipótesis sobre $S$, si $f:X\longrightarrow S$ es normal,
como veremos en (3.5). Consignemos a este respecto:

\result{PROPOSICIÓN}{2.10}

\noindent(i) Si $\operatorname{Pic}_{X/S}$ es liso (respectivamente, plano)
sobre $S$ en los puntos de la sección neutra, también lo es en todo punto de
$\operatorname{Pic}_{X/S}^\sigma$ y en los puntos de toda sección de
$\operatorname{Pic}_{X/S}$ sobre $S$.

\src{236}
\noindent(ii) Sea $s\in S$ tal que
\[
H^2(X_s,\mathcal O_{X_s})=0.
\]
Entonces existe un entorno abierto $U$ de $s$ tal que
$\operatorname{Pic}_{X/S}|U$ es liso sobre $U$.

\noindent(iii) Si $X$ es un esquema propio sobre un cuerpo, entonces
\[
\dim\operatorname{Pic}_{X/k}\leq\dim H^1(X,\mathcal O_X),
\]
con igualdad si y solo si $\operatorname{Pic}_{X/k}$ es liso sobre $k$; este
es siempre el caso cuando $k$ tiene característica cero.

La parte (i) se deduce de (2.5) y (1.12), y (ii) se deduce del criterio
infinitesimal de lisura y de un cálculo conocido de obstrucciones, pues el
«teorema de semicontinuidad» muestra que la hipótesis en $s$ sigue siendo
válida en los puntos próximos. Por último, (iii) se deduce de que
$H^1(X,\mathcal O_X)$ es isomorfo al espacio tangente de Zariski de
$\operatorname{Pic}_{X/k}$ en el elemento neutro. La última afirmación es un
caso particular de un teorema de Cartier que dice que un «grupo formal» en
característica cero es formalmente liso sobre $k$.

\fsection{3}{El subesquema abeliano canónico de
$\operatorname{Pic}_{X/S}$ y el esquema de Albanese}

\result{PROPOSICIÓN}{3.1} Sea $k$ un cuerpo y sea $G$ un esquema conmutativo en
grupos de tipo finito sobre $k$, sin componente aditiva. Entonces
$G_{\mathrm{red}}^\circ$ es separable sobre $k$ y, por consiguiente, es un
subesquema en grupos liso sobre $k$.

Esto es trivial cuando $k$ es perfecto, en particular para $G_{\bar k}$, donde
$\bar k$ es una clausura algebraica de $k$. Basta, por tanto, ver que
$(G_{\bar k}^\circ)_{\mathrm{red}}$ desciende de un subesquema de $G$. De la
hipótesis sobre $G_{\bar k}$ (ausencia de componente aditiva)\footnote{El texto
impreso contiene una construcción mal formada equivalente a «La hipótesis ...
resulta». La fe de erratas vinculada, p.~303, inserta la preposición necesaria
para «De la hipótesis ... resulta».} se deduce fácilmente que existe un entero
$m$ tal que $(G_{\bar k}^\circ)_{\mathrm{red}}$ es la imagen esquemática del
homomorfismo de potencia $m$-ésima de $G_{\bar k}$. Como este homomorfismo
procede del homomorfismo análogo de $G$, la imagen esquemática de este último
da el subesquema buscado.

\result{COROLARIO}{3.2} Sea $X$ un esquema geométricamente normal propio sobre $k$ tal que
exista $\operatorname{Pic}_{X/k}$. Entonces existe un subesquema abeliano $A$
de $\operatorname{Pic}_{X/k}$ cuyo conjunto subyacente es
$\operatorname{Pic}_{X/k}^\circ$.

En efecto, por (2.1), (ii), $\operatorname{Pic}_{X/k}^\circ$ es propio sobre
$k$ y, por tanto, satisface la condición de (3.1).

\src{237}
Así, el resultado anterior muestra que, en ciertos casos, la «variedad de
Picard» clásica (que en la presente teoría se denotaría
$(\operatorname{Pic}_{X/k})_{\mathrm{red}}^\circ$) «está definida sobre $k$»,
sin suponer que $k$ sea perfecto.

Sea ahora $f:X\longrightarrow S$ un esquema relativo propio y plano, con
\[
\mathcal O_S\xrightarrow{\sim}f_*(\mathcal O_X)
\]
para simplificar, tal que exista $\operatorname{Pic}_{X/S}$ y
$\operatorname{Pic}_{X/S}^\circ$ sea propio sobre $S$. Para (3.3), (ii),
suponemos además que existe un abierto de $\operatorname{Pic}_{X/S}$ que
contiene $\operatorname{Pic}_{X/S}^\circ$ y es cuasiproyectivo sobre $S$. Como
se vio anteriormente, esta condición se verifica cuando $f$ es proyectivo con
fibras geométricamente separables e irreducibles. Recuérdese que un esquema
abeliano sobre $S$ es un esquema en grupos sobre $S$ propio y liso sobre $S$,
con fibras geométricamente conexas. Preguntamos si existe un subesquema en
grupos $A$ de $\operatorname{Pic}_{X/S}$ que sea un esquema abeliano y cuyo
conjunto subyacente sea $\operatorname{Pic}_{X/S}^\circ$. Acabamos de ver que
siempre existe uno cuando $S$ es el espectro de un cuerpo. He aquí lo que puede
decirse en el caso general considerado:

\result{TEOREMA}{3.3} Bajo las hipótesis anteriores:

\noindent(i) Si existe un subesquema abeliano de
$\operatorname{Pic}_{X/S}$ cuyo conjunto subyacente sea
$\operatorname{Pic}_{X/S}^\circ$, es único. Por consiguiente, su formación es
compatible con el cambio de base.

\noindent(ii) Para que exista tal subesquema abeliano, basta que exista después
de todo cambio de base $S'\longrightarrow S$ con $S'$ artiniano local. Si $S$
es el espectro de un anillo local, basta comprobar los
$S'=\operatorname{Spec}(A_n)$, donde $A_n=A/\mathfrak m^{n+1}$. Si $S$ es
reducido, basta asimismo comprobar los cambios de base
$S'\longrightarrow S$ en los que $S'$ es el espectro de un anillo de
valoración discreta (completo y con cuerpo residual algebraicamente cerrado,
si se desea).

\noindent(iii) Supongamos que existe $A$, y sea
$B=\operatorname{Alb}^0(X/S)$ el esquema abeliano dual (es decir,
$B=\operatorname{Pic}_{A/S}^\circ$ [7]). Se puede construir entonces,
canónicamente, un espacio homogéneo principal
$P=\operatorname{Alb}^1(X/S)$ bajo $B$, junto con un $S$-morfismo
$X\longrightarrow P$ universal para los $S$-morfismos de $X$ a esquemas
para-abelianos (es decir, espacios homogéneos principales bajo esquemas
abelianos). La formación de $\operatorname{Alb}^0(X/S)$,
$\operatorname{Alb}^1(X/S)$ y
$X\longrightarrow\operatorname{Alb}^1(X/S)$ conmuta con el cambio de base.

Esbocemos la demostración.

La parte (i) es una propiedad general de rigidez de los subesquemas abelianos
de los esquemas conmutativos en grupos: si dos de esos subesquemas coinciden
como conjuntos en un punto $s\in S$, coinciden sobre toda la componente conexa

\src{238}
de $s$ [7]. (Este resultado generaliza un teorema clásico de Chow.)

(ii) Los esquemas de Hilbert muestran que el funtor que asocia a cada $S'$ sobre
$S$ el conjunto (con uno o cero elementos) de los subesquemas abelianos
canónicos de $(\operatorname{Pic}_{X/S})\times_S S'$ está representado por un
esquema $T$ de tipo finito sobre $S$. Por (i), $T\longrightarrow S$ es un
monomorfismo y, por (3.2), es sobreyectivo. La existencia de un subesquema
abeliano canónico de $\operatorname{Pic}_{X/S}$ significa que $T$ posee una
sección sobre $S$, o equivalentemente que $T\longrightarrow S$ es un
isomorfismo. Esto equivale también a que $T\longrightarrow S$ sea étale o,
cuando $S$ es reducido, a que $T\longrightarrow S$ sea propio. Esto demuestra
(ii).

(iii) Directamente de la definición de $\operatorname{Pic}_{X/S}$, para todo
esquema abeliano $C$ sobre $S$, dar un $S$-morfismo de $X$ a un espacio
homogéneo principal bajo $C$ equivale a dar un homomorfismo de grupos
$C'\longrightarrow\operatorname{Pic}_{X/S}$, donde $C'$ es el esquema abeliano
dual de $C$. Si existe el subesquema abeliano canónico $A$ de
$\operatorname{Pic}_{X/S}$, estos homomorfismos se factorizan necesariamente a
través de $A$ (como puede verse también utilizando los puntos de orden finito).
Esto demuestra (iii).

\result{OBSERVACIONES}{3.4} Denotaremos por
$\operatorname{Pic}_{X/S}^{\circ\circ}$ el subesquema abeliano canónico de
$\operatorname{Pic}_{X/S}$, cuando exista. Por desgracia, no siempre existe,
como se ve haciendo variar infinitesimalmente la superficie de Igusa (mediante
la deformación modular de primer orden). Es posible, no obstante, que
$\operatorname{Pic}_{X/S}^{\circ\circ}$ exista al menos cuando $S$ sea reducido
o, equivalentemente por (ii), cuando $S$ sea el espectro de un anillo de
valoración discreta. Sean $X_0,X_1$ las fibras especial y genérica de $X$, y
pongamos
\[
A_1=\operatorname{Pic}_{X_1/K}^\circ,
\]
donde $K$ es el cuerpo de fracciones del anillo de valoración $V$. Por el
teorema de Koizumi, existe un esquema abeliano $A$ sobre $S$, esencialmente
único, cuya fibra genérica es $A_1$. Como en (2.1), (i), se ve fácilmente
(suponiendo ahora que $X$ es liso sobre $S$) que el morfismo identidad de $A_1$
se prolonga a un morfismo
\[
A\longrightarrow\operatorname{Pic}_{X/S}.
\]
De ahí se obtiene un homomorfismo
\[
\tag{*}A_0\longrightarrow\operatorname{Pic}_{X_0/k}^{\circ\circ},
\]
que se demuestra fácilmente que es sobreyectivo y tiene por núcleo un
$p$-grupo finito, donde $p$ es el exponente característico del cuerpo residual
$k$

\src{239}
(de nuevo utilizando los puntos de orden finito). Sentado esto, las condiciones
siguientes sobre $X/S$ son equivalentes:

\medskip
\noindent a. El homomorfismo anterior (*) es un isomorfismo (lo que Shimura
expresaría diciendo que la formación de la «variedad de Picard» es «compatible
con la especialización»).

\noindent b. Existe $\operatorname{Pic}_{X/S}^{\circ\circ}$ (en cuyo caso no es
otro que $A$).

\noindent c. (Para referencia) Existen los
$\operatorname{Pic}_{X_n/S}^{\circ\circ}$.

Por la observación hecha sobre el núcleo de (*), la condición (a) se verifica
cuando la característica residual es cero, pero este resultado se generalizará
considerablemente en (3.5).

Por supuesto, si $\operatorname{Pic}_{X/S}$ es liso sobre $S$ en los puntos de
$\operatorname{Pic}_{X/S}^\circ$, este último es abierto en
$\operatorname{Pic}_{X/S}$ (cf. (1.7)). Provisto de la estructura inducida, es
por tanto un subesquema abeliano de $\operatorname{Pic}_{X/S}$, luego igual a
$\operatorname{Pic}_{X/S}^{\circ\circ}$, que existe en este caso. Pero se tiene
algo mucho más fuerte:

\result{TEOREMA}{3.5} Bajo las hipótesis de (3.3), sea $s\in S$ tal que
$\operatorname{Pic}_{X_s/k(s)}$ sea liso sobre $k(s)$ (o, equivalentemente,
\[
\dim\operatorname{Pic}_{X_s/k(s)}=\dim H^1(X_s,\mathcal O_{X_s})).
\]
Entonces existe un entorno abierto $U$ de $s$ tal que
$\operatorname{Pic}_{X/S}$ es liso sobre $S$ en los puntos de
$\operatorname{Pic}_{X/S}^\circ|U$, que es, por consiguiente, un subesquema
abeliano abierto de $\operatorname{Pic}_{X/S}|U$. A fortiori, existe
$\operatorname{Pic}_{X|U/U}^{\circ\circ}$.

Indiquemos el principio de la demostración. La discusión anterior nos reduce
al caso en que $S$ es el espectro de un anillo artiniano local $A$, y se razona
por inducción sobre el orden infinitesimal de $A$. Podemos suponer, por tanto,
que $\operatorname{Pic}_{X_n/A_n}^\circ$ es liso sobre $A_n$ y basta demostrar
que $\operatorname{Pic}_{X_{n+1}/A_{n+1}}^\circ$ es liso sobre $A_{n+1}$. Para
ello basta construir un esquema abeliano $P_{n+1}$ sobre $A_{n+1}$ que prolongue
\[
P_n=\operatorname{Pic}_{X_n/A_n}^\circ,
\]
junto con un Módulo invertible $\mathcal L_{n+1}$ sobre
$X_{n+1}\times_{A_{n+1}}P_{n+1}$ que prolongue el Módulo invertible
$\mathcal L_n$ sobre $X_n\times_{A_n}P_n$ que interviene en la definición del
esquema de Picard $\operatorname{Pic}_{X_n/A_n}$ como solución de un problema
universal. (N.B. --- Podemos suponer que $X$ posee una sección sobre $S$~...)
La construcción utiliza el resultado clave siguiente: todo esquema abeliano
definido sobre un cociente de un anillo artiniano local puede prolongarse (en
otros términos, el «esquema formal de Módulos» absoluto ([3], II) de un esquema
abeliano sobre un cuerpo algebraicamente cerrado es liso sobre el anillo de

\src{240}
vectores de Witt de $k$). Este resultado se deduce, a su vez, de las propiedades
formales generales de la obstrucción al levantamiento y de las operaciones de
grupo. Admitiéndolo, se prolonga primero $P_n$ arbitrariamente a $P_{n+1}$. Se
encuentra entonces una obstrucción al levantamiento de $\mathcal L_n$, que
pertenece a
\[
H^2(X_0\times P_0,\mathcal O_{X_0\times P_0})\otimes_k V
\qquad(\text{donde }V=\mathfrak m^{n+1}/\mathfrak m^{n+2}),
\]
y, de hecho, al subespacio
\[
H^1(X_0,\mathcal O_{X_0})\otimes
H^1(P_0,\mathcal O_{P_0})\otimes_k V
\]
(porque la restricción de $\mathcal L_n$ a cada uno de los dos factores $X_n$
y $P_n$ es trivial). Este último espacio es precisamente
\[
H^1(P_0,\mathcal T_{P_0/k})\otimes V,
\]
donde $\mathcal T_{P_0/k}$ es el haz tangente de $P_0/k$, y es también, por
tanto, el espacio que mide la indeterminación en el levantamiento de $P_n$ a
$P_{n+1}$ ([3], II). Por consiguiente, este levantamiento puede corregirse (de
una y solo una manera, como debe ser) para anular la obstrucción al
levantamiento de $\mathcal L_n$.

\result{COROLARIO}{3.6} Bajo las hipótesis de (3.5),
$R^1f_*(\mathcal O_X)$ es un Módulo localmente libre sobre $S$ en un entorno de
$s$, y su formación conmuta con el cambio de base.

En efecto, este Módulo es el Módulo tangente de $\operatorname{Pic}_{X/S}$ a lo largo
de la sección neutra.

\result{OBSERVACIÓN}{3.7} El mismo argumento que en (3.5) muestra que, si $S'$
es un subesquema de $S$ definido por un Ideal coherente nilpotente y se supone
solamente que existe $\operatorname{Pic}_{X'/S'}^\circ$ y que tiene fibras
lisas, entonces existe necesariamente $\operatorname{Pic}_{X/S}^\circ$ y es un
esquema abeliano sobre $S$. Esto permite construir esquemas de Picard en
algunos casos pese a la ausencia de una hipótesis de proyectividad; por
ejemplo, siempre existe el esquema abeliano dual de un esquema abeliano sobre
un anillo artiniano. Por otra parte, aplicando (3.6) cuando $X$ es un esquema
abeliano sobre $S$ y utilizando la estructura conocida de álgebra exterior de
$H^*(X_s,\mathcal O_{X_s})$ sobre $H^1(X_s,\mathcal O_{X_s})$
(Rosenlicht--Serre), se obtiene que, para todo $i$, $R^if_*(\mathcal O_X)$ es
localmente libre y, más precisamente, isomorfo a la potencia exterior $i$-ésima
de $R^1f_*(\mathcal O_X)$.

\result{OBSERVACIÓN}{3.8} Para un morfismo proyectivo liso
$f:X\longrightarrow S$, con $S$ reducido y todas las características
residuales nulas, (3.6) era conocido por métodos trascendentes como consecuencia
de la teoría de Hodge. De hecho, todos los $R^pf_*(\Omega^q_{X/S})$ son
entonces localmente libres. Por otra parte, existen contraejemplos

\src{241}
en característica mixta a la libertad local de $R^1f_*(\mathcal O_X)$,
obtenidos a partir de variedades de Serre ([8], p.~50). No parece haber
contraejemplos en característica igual.

\result{COROLARIO}{3.8} Bajo las hipótesis de (3.5),
$\operatorname{Pic}_{X/S}|U$ es liso sobre $U$ en todo punto de
$\operatorname{Pic}_{X/S}^\sigma|U$.

Aplíquese (2.10), (i). En particular, teniendo en cuenta (2.10), (iii):

\result{COROLARIO}{3.9} Supongamos que todas las características residuales de
$S$ son nulas. Entonces $\operatorname{Pic}_{X/S}^\tau$ es liso sobre $S$.

\result{OBSERVACIÓN}{3.10} Si $\operatorname{Pic}_{X/S}^\tau$ es además propio sobre
$S$, se sigue, por ejemplo, que el subgrupo de torsión de los grupos de
Néron--Severi de las fibras geométricas de $f$ es constante sobre cada
componente conexa de $S$ (como también resulta claro por métodos trascendentes
cuando $f$ es liso y proyectivo). Una aplicación directa de (2.5) muestra, más
generalmente, que si $\operatorname{Pic}_{X/S}^\tau$ es propio sobre $S$ (por
ejemplo, si $f:X\longrightarrow S$ es liso y proyectivo) y $q$ es un primo
distinto de todas las características residuales de $S$, entonces la
componente $q$-primaria de los subgrupos de torsión de los grupos de
Néron--Severi de las fibras geométricas de $X/S$ es constante sobre cada
componente conexa de $S$. En característica $p>0$, esto deja de ser cierto para
la componente $p$-primaria de la torsión. Sigue siendo posible, sin embargo,
que el rango sobre $k(s)$ de
\[
\operatorname{Pic}_{X_s/k(s)}^\tau/
\operatorname{Pic}_{X_s/k(s)}^{\circ\circ}=T_{X_s/k(s)}
\]
sea localmente constante. Cuando $S$ es reducido, se puede demostrar que esto
equivale a la existencia de $\operatorname{Pic}_{X/S}^{\circ\circ}$ junto con
la planitud de $\operatorname{Pic}_{X/S}^\tau$ sobre $S$, y basta comprobar la
cuestión cuando $S$ es el espectro de un anillo de valoración discreta. Esto es
lo que he comprobado en los pocos ejemplos que he examinado; pero como el
enunciado correspondiente con $S$ artiniano es falso (cf. Observaciones (2.9)
y (3.4)), probablemente no convenga entregarse a un optimismo excesivo.

\fsection{4}{El teorema de finitud para el esquema de Picard}

Sea $f:X\longrightarrow S$ un morfismo proyectivo y plano tal que exista
$\operatorname{Pic}_{X/S}$. La consideración de los «polinomios de Hilbert»
$Q\in\mathbf Q[t]$ descompone $\operatorname{Pic}_{X/S}$ en una unión disjunta
de abiertos $\operatorname{Pic}_{X/S}^Q$. Sin hipótesis adicionales, por
ejemplo que aseguren que $\operatorname{Pic}_{X/S}$ sea separado sobre $S$,
estos abiertos no son necesariamente de tipo finito sobre $S$; se obtiene un
contraejemplo cuando $X$ es «una cónica que degenera en dos rectas». Es

\src{242}
posible, sin embargo, que sean de tipo finito cuando $f$ sea separable y tenga
fibras geométricamente irreducibles. La cuestión está relacionada con saber si
$\operatorname{Pic}_{X/S}^\tau$ es de tipo finito sobre $S$, lo que podría ser
cierto sin ninguna hipótesis sobre las fibras de $X/S$. Cuando
$f:X\longrightarrow S$ es liso, obsérvese que
$\operatorname{Pic}_{X/S}^\tau$ está contenido en uno de los
$\operatorname{Pic}_{X/S}^Q$ (pues sobre una variedad proyectiva no singular la
equivalencia por torsión es más fina que la equivalencia numérica para los
divisores y, de hecho, por Matsusaka, coincide con ella), y es, por tanto, de
tipo finito sobre $S$ si los $\operatorname{Pic}_{X/S}^Q$ lo son. Las
cuestiones de finitud que acabamos de plantear conservan un sentido evidente
incluso sin suponer la existencia de $\operatorname{Pic}_{X/S}$: dicen que
ciertas familias de Módulos invertibles son «acotadas» en el sentido de [3], IV.
Para todo cuerpo algebraicamente cerrado $k$, considérense los subesquemas
íntegros de $\mathbf P_k^r$ de dimensión $n$ y grado $d$, y los Módulos
invertibles sobre esos preesquemas con polinomio de Hilbert $Q$, donde
$r,n,d,Q$ están fijados. El problema consiste en demostrar que la familia de
estos Módulos (considerados como Módulos coherentes sobre los $\mathbf P_k^r$) es
acotada, es decir, que puede parametrizarse mediante un esquema de tipo finito
sobre $\mathbf Z$.

Utilizando el método de Matsusaka [6], una demostración bastante técnica (que
emplea los «criterios de equivalencia» en una forma adecuada) da una respuesta
afirmativa cuando uno se restringe a las subvariedades no singulares de
$\mathbf P_k^r$. Más precisamente:

\result{TEOREMA}{4.1} Sea $f:X\longrightarrow S$\footnote{El texto impreso
dice «Sea: $X\longrightarrow S$»; el morfismo utilizado más adelante es $f$.}
un morfismo proyectivo liso con fibras geométricamente conexas, siendo $S$
noetheriano. Sea $\mathcal O_X(1)$ un Módulo relativamente muy amplio sobre $X$,
y sea $E$ un subconjunto de $\operatorname{Pic}_{X/S}$ correspondiente a una
familia $(\mathcal L_i)$ de Módulos invertibles sobre las fibras geométricas
$\overline X_{s_i}$ de $X/S$. Sea $D_i$ un divisor (no necesariamente
efectivo) sobre $\overline X_{s_i}$ que defina $\mathcal L_i$, sea
\[
a_n^{(i)}t^n+\cdots+a_0^{(i)}
\]
el polinomio de Hilbert de $\mathcal L_i$, y sea $\xi=\xi_i$ un divisor que
defina $\mathcal O_X(1)$, es decir, una sección hiperplana. Las condiciones
siguientes son equivalentes:

\noindent a. $E$\footnote{El texto impreso dice «$Q$». El subconjunto
introducido en el enunciado, cuya cuasicompacidad expresa la acotación de la
familia, es $E$.} es cuasicompacto, es decir, está contenido en un abierto de
tipo finito sobre $S$; equivalentemente, la familia $(\mathcal L_i)$ es
acotada.

\noindent b. $E$ está contenido en la unión de un número finito de los
conjuntos $\operatorname{Pic}_{X/S}^Q$, $Q\in\mathbf Q[t]$; equivalentemente,
solo aparece un número finito de polinomios de Hilbert entre los
$\mathcal L_i$.

\noindent b'. (Si todas las fibras de $X/S$ tienen la misma dimensión $n$.)
Los dos coeficientes $a_{n-1}^{(i)}$ y $a_{n-2}^{(i)}$ de los polinomios de
Hilbert de los $\mathcal L_i$ recorren conjuntos finitos.

\src{243}
\noindent b''. (Si todas las fibras de $X/S$ tienen la misma dimensión $n$.)
Los enteros
\[
\xi^{n-1}D_i\qquad\text{y}\qquad\xi^{n-2}D_i^2
\]
(los grados proyectivos de $D_i$ y $D_i^2$) permanecen acotados superiormente.

\result{COROLARIO}{4.2} Sea $f:X\longrightarrow S$ un morfismo proyectivo liso
con fibras geométricamente conexas. Entonces los esquemas
$\operatorname{Pic}_{X/S}^Q$ y $\operatorname{Pic}_{X/S}^\tau$
($Q\in\mathbf Q[t]$) son proyectivos sobre $S$.

Son de tipo finito sobre $S$ por (4.1), por lo que se aplica (2.4).

\section*{BIBLIOGRAFÍA}
\addcontentsline{toc}{section}{Bibliografía}

\begin{description}
\item[{[1]}] CHEVALLEY (Claude). -- Sur la théorie de Picard,
\emph{Amer. J. of Math.}, vol. 82, 1960, pp. 435--490.

\item[{[2]}] GROTHENDIECK (A.) y DIEUDONNÉ (J.). -- Éléments de géométrie
algébrique, capítulo I y siguientes. -- París, Presses universitaires de
France, 1960--1961 (Institut des Hautes Études Scientifiques, 4, 8, 11,
\ldots).

\item[{[3]}] GROTHENDIECK (Alexander). -- Technique de descente et théorèmes
d'existence en géométrie algébrique, I--V, \emph{Séminaire Bourbaki}, t. 12,
1959/60, n.º 190, 29 pp., y n.º 195, 22 pp.; t. 13, 1960/61, n.º 212,
20 pp., y n.º 221, 28 pp.; t. 14, 1961/62, n.º 232, 19 pp.

\item[{[4]}] GROTHENDIECK (Alexander). -- Séminaire de géométrie algébrique,
1960/61. -- París, Institut des Hautes Études Scientifiques, 1961.

\item[{[5]}] GROTHENDIECK (Alexander). -- Séminaire de géométrie algébrique,
1961/62, redactado por un grupo de oyentes. -- París, Institut des Hautes Études
Scientifiques (en preparación).

\item[{[6]}] MATSUSAKA (T.). -- The criteria for algebraic equivalence and the
torsion group, \emph{Amer. J. Math.}, vol. 79, 1957, pp. 53--66.

\item[{[7]}] MUMFORD (D.) y TATE (J.). -- Séminaire de géométrie algébrique,
Harvard University, semestre de primavera de 1962 (en preparación).

\item[{[8]}] SERRE (Jean-Pierre). -- Sur la topologie des variétés algébriques
en caractéristique $p$, \emph{Symposium internacional de topología algebraica}
[1956, México], pp. 24--53. -- México, Universidad Nacional Autónoma, 1958.
\end{description}
